% AUTO-GENERATED by corpus/make_appendix_table.py — do not edit by hand.
\begingroup
\renewcommand{\thetable}{A.\arabic{table}}
\setcounter{table}{0}
\begin{landscape}
\scriptsize
\begin{longtable}{@{}l p{2.6cm} p{3.4cm} p{1.8cm} p{5.0cm} p{2.1cm} p{6.2cm}@{}}
\caption{Included studies: characteristics and design (n = 67). Source: author's compilation.}\label{tab:coding-a1}\\
\toprule
ID & Study & Journal (AJG) & Country & Sample & Method & AI investment measure \\
\midrule
\endfirsthead
\multicolumn{7}{@{}l}{\tablename~\thetable{} (continued)}\\
\toprule
ID & Study & Journal (AJG) & Country & Sample & Method & AI investment measure \\
\midrule
\endhead
\bottomrule
\endfoot
S01 & Wamba-Taguimdje et al. (2020) & Business Process Management Journal (2) & multi-country & \textasciitilde{}500 vendor-published AI project cases & case study & case observation (vendor project reports) \\
\addlinespace[2pt]
S02 & Baabdullah et al. (2021) & Industrial Marketing Management (3) & Saudi Arabia & 392 B2B SMEs & survey-SEM & survey construct (acceptance of AI practices) \\
\addlinespace[2pt]
S03 & Chatterjee et al. (2021) & Industrial Marketing Management (3) & India & 349 managers & survey-SEM & survey construct (AI-CRM implementation) \\
\addlinespace[2pt]
S04 & Chatterjee et al. (2022) & Journal of Business Research (3) & India & 312 manager responses & survey-SEM & survey construct (AI-embedded CRM adoption) \\
\addlinespace[2pt]
S05 & Hossain et al. (2022) & Industrial Marketing Management (3) & Bangladesh & 257 manager responses & survey-SEM & survey construct (AI adoption on a marketing analytics platform) \\
\addlinespace[2pt]
S06 & Lee et al. (2022) & Technovation (3) & South Korea & 160 high-tech firms & panel econometrics & survey (AI-adoption intensity) linked to financial records \\
\addlinespace[2pt]
S07 & Leoni et al. (2022) & International Journal of Operations and Production Management (4) & Italy & 120 senior executives & survey-SEM & survey construct (AI adoption/use) \\
\addlinespace[2pt]
S08 & Lui et al. (2022) & Annals of Operations Research (3) & USA & 119 AI investment announcements by 62 listed firms, 2015-2019 & event study & announcements (AI investment announcements) \\
\addlinespace[2pt]
S09 & Mishra et al. (2022) & Journal of the Academy of Marketing Science (4*) & USA & 19,000 firm-year observations of US-listed firms & panel econometrics & annual-report text (share of AI-related keywords) \\
\addlinespace[2pt]
S10 & Sun et al. (2022) & Systems Research and Behavioral Science (2) & China & 234 AI-related manufacturers (general-manager/CEO respondents) & survey-SEM & survey construct (value co-creation in AI innovation ecosystem) \\
\addlinespace[2pt]
S11 & Ali et al. (2023) & Journal of Organizational Change Management (2) & UAE & single hospital case (Dubai): 9 interviews with the robotic-surgery team plus archival data & case study & case observation (adoption of AI/robotic surgery) \\
\addlinespace[2pt]
S12 & Czarnitzki et al. (2023) & Journal of Economic Behavior and Organization (3) & Germany & 5,851 firms (409 AI users), representative German innovation survey 2018 & panel econometrics & survey (AI adoption yes/no plus intensity of AI use in business processes) \\
\addlinespace[2pt]
S13 & Huang \& Lee (2023) & Data Base for Advances in Information Systems (2) & USA & 109 AI implementation announcements by S\&P 500 firms, 2014-2019 & event study & announcements (AI implementation) \\
\addlinespace[2pt]
S14 & Krakowski et al. (2023) & Strategic Management Journal (4*) & multi-country & same chess players across human-only, human-plus-AI and engine-only tournaments (controlled competitive setting) & panel econometrics & adoption of chess engines (AI) in competitive play \\
\addlinespace[2pt]
S15 & Samadhiya et al. (2023) & Industrial Marketing Management (3) & India & interviews with senior B2B managers plus a survey of 142 responses & mixed & survey construct (AIPRM capabilities: AI-based partner relationship management) \\
\addlinespace[2pt]
S16 & Wu \& Monfort (2023) & Psychology and Marketing (3) & Taiwan & 278 CEO/marketing-manager responses, food franchises and chain stores, 2020 & survey-SEM & survey construct (implementation of AI marketing strategy) \\
\addlinespace[2pt]
S17 & Babina et al. (2024) & Journal of Financial Economics (4*) & USA & 1,993 US listed firms & panel econometrics & resume-based (share of AI-skilled workers as a proxy for firm AI investment) \\
\addlinespace[2pt]
S18 & Cannas et al. (2024) & International Journal of Production Research (3) & Italy & 6 companies, 17 AI implementation cases, interview-based & case study & case observation (AI applications across supply-chain processes) \\
\addlinespace[2pt]
S19 & Pesqueira et al. (2024) & Computers and Industrial Engineering (2) & Europe & 4 pharmaceutical companies (2 AI adopters vs. 2 non-adopters) plus a two-round expert panel of 53 & mixed & adoption status (AI in tender management) \\
\addlinespace[2pt]
S20 & Pham et al. (2024) & Journal of Business Research (3) & USA & 941 AI and 941 matched non-AI hospital-year observations (246 hospitals), US, 2000-2020 & panel econometrics & adoption dummy (hospital has adopted AI or not) \\
\addlinespace[2pt]
S21 & Sullivan \& Fosso (2024) & Journal of Business Research (3) & Europe & 107 executives (France + UK), two-wave survey & survey-SEM & survey constructs (AI-enabled automation, AI-enabled analytics, AI-enabled relational capabilities) \\
\addlinespace[2pt]
S22 & Sun et al. (2024) & Technology in Society (2) & China & 3,235 Chinese listed firms, 2007-2021 & panel econometrics & annual-report text (frequency of AI-related terms) \\
\addlinespace[2pt]
S23 & Zebec \& Indihar (2024) & Business Process Management Journal (2) & EU & 448 organisations, EU & survey-SEM & survey construct (AI adoption) \\
\addlinespace[2pt]
S24 & Alam et al. (2025) & International Journal of Hospitality Management (3) & Malaysia & 336 hospitality-industry respondents & survey-SEM & survey construct (AI adoption + technology readiness) \\
\addlinespace[2pt]
S25 & Alwakid \& Dahri (2025) & Technology in Society (2) & Saudi Arabia & 250 SMEs (managers/owners), manufacturing + services & survey-SEM & survey construct (AI capabilities: infrastructure, business integration, proactive stance) \\
\addlinespace[2pt]
S26 & Banna \& Alam (2025) & International Journal of Entrepreneurial Behaviour and Research (3) & EU & 1,479 firms in 26 EU countries, 2012-2023 & panel econometrics & AI venture capital invested in the firm's country (same value for every firm in that country) \\
\addlinespace[2pt]
S27 & Basnet et al. (2025) & International Review of Financial Analysis (3) & USA & US firms, annual reports (10-K) 2005–2018 & panel econometrics & annual-report text (AI mentions classified as actionable, speculative, or irrelevant) \\
\addlinespace[2pt]
S28 & Bin-Nashwan \& Li (2025) & Technology in Society (2) & China & 433 Chinese accounting professionals, survey 2023–2024 & survey-SEM & survey construct (AI-infused knowledge systems) \\
\addlinespace[2pt]
S29 & Cao et al. (2025) & European Management Review (3) & USA & 140 US retail managers; instrument built from 6,519 AI documents of 37 retailers & fsQCA & survey instrument (AI integration at task, function, and firm level) \\
\addlinespace[2pt]
S30 & Chakraborty (2025) & Current Issues in Tourism (2) & India & 408 hotel and resort managers, India, two-wave survey & survey-SEM & survey construct (GenAI adoption) \\
\addlinespace[2pt]
S31 & Chiu et al. (2025) & Business Ethics the Environment and Responsibility (2) & Taiwan & 33 listed semiconductor firms & DEA & patents (AI-driven innovation proxied by patent activity) \\
\addlinespace[2pt]
S32 & D’Amico et al. (2025) & Journal of Technology Transfer (3) & UK & 14,143 UK firms, 2004-2020 & panel econometrics & archival records (AI technology or product use identified from a commercial database) \\
\addlinespace[2pt]
S33 & Feng et al. (2025) & Quarterly Review of Economics and Finance (2) & China & 4,004 Chinese listed firms, 2011–2023 & panel econometrics & procurement index (AI adoption from procurement records) \\
\addlinespace[2pt]
S34 & Fontanelli et al. (2025) & Journal of Economic Behavior and Organization (3) & France & French firms, 2019 survey on information and communication technology use & panel econometrics & survey (predictive AI use; buyers vs in-house developers) \\
\addlinespace[2pt]
S35 & Guo et al. (2025) & International Marketing Review (3) & China & 322 SMEs using cross-border platforms & survey-SEM & survey construct (AI adoption coupled with strategic agility) \\
\addlinespace[2pt]
S36 & Huang \& Lin (2025) & Managerial and Decision Economics (2) & USA & 50 S\&P 500 AI adopters plus non-adopter controls, 2010–2017 & panel econometrics & news announcements (manually verified AI implementation) \\
\addlinespace[2pt]
S37 & Kumar et al. (2025) & Journal of Business Research (3) & multi-country & 277 B2B managers, multi-country survey & survey-SEM & survey construct (GenAI adoption) \\
\addlinespace[2pt]
S38 & Mehta et al. (2025) & Journal of Business and Industrial Marketing (2) & India & 496 survey respondents plus 26 key-account-manager interviews, automobile sector & mixed & survey construct (adoption of AI technologies by key account managers) \\
\addlinespace[2pt]
S39 & N. \& Roy (2025) & European Journal of Marketing (3) & Australia & 335 companies across industries & survey-SEM & survey construct (AI-based stakeholder engagement) \\
\addlinespace[2pt]
S40 & Sandeep et al. (2025) & Journal of Intellectual Capital (2) & India & 290 HR professionals in talent acquisition, South India & survey-SEM & survey construct (AI adoption in recruitment) \\
\addlinespace[2pt]
S41 & Shi et al. (2025) & International Review of Financial Analysis (3) & China & Chinese listed companies, 2010–2022 (\textasciitilde{}22,900 firm-years) & panel econometrics & annual-report text (dictionary of AI-related terms) \\
\addlinespace[2pt]
S42 & Song et al. (2025) & Industrial Management and Data Systems (2) & USA & 120 AI service-failure events at US-listed companies, 2010–2024 & event study & announcements (AI service-failure events) \\
\addlinespace[2pt]
S43 & Tao et al. (2025) & Journal of Business and Industrial Marketing (2) & China & 291 B2B manufacturing firms, multi-wave, multi-source survey & survey-SEM & survey construct (big-data analytics assimilation and AI capabilities) \\
\addlinespace[2pt]
S44 & Tingbani et al. (2025) & International Journal of Finance and Economics (3) & USA & 1,950 US firms, 1996–2016 & panel econometrics & AI hiring intensity: share of AI skills in the firm's workforce, from job postings and employee-profile (resume) data \\
\addlinespace[2pt]
S45 & Adams et al. (2026) & Journal of Monetary Economics (4) & USA & US firms matched to job postings, 2010–2024 & panel econometrics & job postings (openings combining AI skills and pricing work signal AI-pricing adoption) \\
\addlinespace[2pt]
S46 & Agag et al. (2026) & Tourism Management (4) & UK & 1,206 UK listed tourism, travel, and hospitality firms, 2018–2024 & panel econometrics & AI human capital (AI-skilled workforce measure) \\
\addlinespace[2pt]
S47 & Alnofeli et al. (2026) & International Journal of Information Management (2) & Australia & 205 banking employees & survey-SEM & survey construct (AI-powered CRM capability) \\
\addlinespace[2pt]
S48 & Arshad et al. (2026) & International Journal of Information Management (2) & multi-country & two cross-sectional surveys of large firms worldwide & survey-based regression (two cross-sectional surveys) & survey construct (joint implementation of robotic process automation and AI) \\
\addlinespace[2pt]
S49 & Bai et al. (2026) & Pacific Basin Finance Journal (2) & China & Chinese listed firms; AI pilot-zone policy as natural experiment & panel econometrics & policy shock (location in a government AI pilot zone as driver of AI investment) \\
\addlinespace[2pt]
S50 & Chen et al. (2026) & Journal of Empirical Finance (3) & USA & US firms & panel econometrics & earnings-call text (share of AI terminology in management remarks) \\
\addlinespace[2pt]
S51 & Dinh (2026) & Industrial Marketing Management (3) & Vietnam & 261 manufacturing SMEs & mixed & survey construct (AI-enabled knowledge capabilities) \\
\addlinespace[2pt]
S52 & Filieri et al. (2026) & Journal of Product Innovation Management (4) & China & Chinese listed manufacturing exporters, 2015-2022 & panel econometrics & archival AI integration measure (adoption induced by tariff exposure) \\
\addlinespace[2pt]
S53 & Giordino et al. (2026) & Technological Forecasting and Social Change (3) & Europe & 432 European listed firms, 2015-2023 (financial firms excluded) & panel econometrics & archival score of organizational AI focus (commercial database rating) \\
\addlinespace[2pt]
S54 & Kazakis (2026) & Scottish Journal of Political Economy (2) & USA & US listed firms & panel econometrics & labor-based measure (AI investment inferred from the firm's AI-related workforce) \\
\addlinespace[2pt]
S55 & Li et al. (2026) & International Journal of Production Economics (3) & China & 218 firms in supply chains & survey-SEM & survey construct (responsible AI adoption) \\
\addlinespace[2pt]
S56 & Li et al. (2026) & Journal of Business Research (3) & China & 223 firms, cross-industry & survey-SEM & survey construct (GenAI affordances: organizational memory, collaboration, creation) \\
\addlinespace[2pt]
S57 & Lin et al. (2026) & Emerging Markets Finance and Trade (2) & China & Chinese listed firms, 2007-2022 & panel econometrics & annual-report text (AI adoption from AI-related wording) \\
\addlinespace[2pt]
S58 & Lin et al. (2026) & International Review of Financial Analysis (3) & China & Chinese listed firms & panel econometrics & annual-report text (AI keyword dictionary) \\
\addlinespace[2pt]
S59 & Liu et al. (2026) & International Review of Financial Analysis (3) & China & Chinese publicly listed firms, 2007-2024 & panel econometrics & annual-report text (frequency of AI-related keywords) \\
\addlinespace[2pt]
S60 & Liu et al. (2026) & Technology in Society (2) & Asia-Pacific & 420 firms (survey) plus a vignette experiment with 360 managers & mixed & survey construct (AI-enabled innovation); experimentally manipulated regulatory conditions \\
\addlinespace[2pt]
S61 & Mukherjee et al. (2026) & International Journal of Quality and Reliability Management (2) & India & 312 responses from luxury hotels & survey-SEM & survey construct (AI adoption intention) \\
\addlinespace[2pt]
S62 & Patel (2026) & Technological Forecasting and Social Change (3) & USA & 1,306 AI-scored service patents from 3,899 manufacturing firms, 1980-2020 & event study & patents (AI intensity of service-technology patents) \\
\addlinespace[2pt]
S63 & Renfei et al. (2026) & Technovation (3) & China & 292 responses from 120 manufacturing firms & survey-SEM & survey construct (perceptive, predictive and prescriptive AI capabilities) \\
\addlinespace[2pt]
S64 & Singh et al. (2026) & Business Strategy and the Environment (3) & USA & S\&P 500 firms, 2000-2024 & panel econometrics & patents (AI patents granted per year) \\
\addlinespace[2pt]
S65 & Song et al. (2026) & Finance Research Letters (2) & USA & US firms, 2018-2023 & panel econometrics & disclosure-vs-patents gap (AI claims not backed by AI patents, i.e. AI washing) \\
\addlinespace[2pt]
S66 & Ullah et al. (2026) & International Journal of Information Management (2) & China & 419 responses from 290 manufacturing firms & survey-SEM & survey construct (AI adoption) \\
\addlinespace[2pt]
S67 & Wang \& Zhang (2026) & International Journal of Information Management (2) & China + Europe & 357 early-stage ventures, two-wave survey & mixed & survey construct (AI innovation capability) \\
\addlinespace[2pt]
\end{longtable}

\clearpage
\begin{longtable}{@{}l p{2.4cm} l p{4.2cm} l p{6.2cm} p{6.2cm}@{}}
\caption{Included studies: outcomes, effect directions, and conditions (n = 67). Source: author's compilation.}\label{tab:coding-a2}\\
\toprule
ID & Study & Outcome & Outcome measure(s) & Direction & Conditions & Key finding \\
\midrule
\endfirsthead
\multicolumn{7}{@{}l}{\tablename~\thetable{} (continued)}\\
\toprule
ID & Study & Outcome & Outcome measure(s) & Direction & Conditions & Key finding \\
\midrule
\endhead
\bottomrule
\endfoot
S01 & Wamba-Taguimdje et al. (2020) & Perf. & Perf: organizational and process performance (qualitative) & conditional & process reconfiguration required; complementary resources (data, talent, domain knowledge, partnerships, infrastructure) & AI transformation projects improve organizational and process performance, but only when firms use AI to reconfigure their processes rather than overlay it on existing ones. \\
\addlinespace[2pt]
S02 & Baabdullah et al. (2021) & Perf. & Perf: perceived SME performance (survey scale) & positive & adoption enablers: technology roadmapping, attitude, infrastructure readiness, awareness; professional expertise not significant & Acceptance of AI practices improves perceived SME performance; adoption is driven by roadmapping, attitude, infrastructure and awareness - not by professional expertise. \\
\addlinespace[2pt]
S03 & Chatterjee et al. (2021) & Both & Perf: perceived organizational performance; CA: competitive advantage scale (incl. market-share gain) & positive & mediators: B2B engagement, employee experience, information processing; moderator: leadership support & AI-CRM implementation improves perceived organizational performance and competitive advantage in B2B relationship management. \\
\addlinespace[2pt]
S04 & Chatterjee et al. (2022) & Perf. & Perf: firm performance scale; B2B relationship satisfaction & positive & technology turbulence weakens the paths to relationship satisfaction; leadership support strengthens the satisfaction-to-performance path & AI-embedded CRM improves relationship satisfaction and firm performance; gains shrink under technology turbulence and grow with leadership support. \\
\addlinespace[2pt]
S05 & Hossain et al. (2022) & CA & CA: sustained competitive advantage (relative performance vs. competitors) & conditional & mediators: market sensing, seizing, and reconfiguring carry the effect in full; moderator: AI adoption strengthens all three capability paths & Marketing analytics capability drives sensing/seizing/reconfiguring toward sustained competitive advantage; AI adoption amplifies this only when built on an analytics foundation. \\
\addlinespace[2pt]
S06 & Lee et al. (2022) & Perf. & Perf: revenue growth & conditional & adoption-intensity threshold (low-level adoption shows no effect); complementary cloud and database investments amplify & AI pays off only above an adoption-intensity threshold, and more so with complementary technology investments and venture-specific internal R\&D. \\
\addlinespace[2pt]
S07 & Leoni et al. (2022) & Perf. & Perf: perceived manufacturing firm performance & conditional & full mediation: knowledge management processes carry the entire effect, the direct path from AI to performance is not significant; AI has no direct effect on supply chain resilience either; larger firms reach higher AI maturity & AI alone does not improve manufacturing firm performance - the effect exists exclusively through knowledge management processes (full mediation). \\
\addlinespace[2pt]
S08 & Lui et al. (2022) & Perf. & Perf: firm market value (abnormal stock returns around the announcement) & negative & stronger negative reaction for nonmanufacturing firms, weak IT capability and low credit ratings; better IT capability lessens the drop & Investors react negatively to AI investment announcements on average; the penalty concentrates in nonmanufacturing firms with weak IT capability and low credit ratings. \\
\addlinespace[2pt]
S09 & Mishra et al. (2022) & Perf. & Perf: gross and net operating efficiency, net profitability, advertising spending, employment & mixed &  & AI focus reduces gross operating efficiency but increases net operating efficiency and profitability - automation raises costs more slowly than revenues. \\
\addlinespace[2pt]
S10 & Sun et al. (2022) & CA & CA: perceived competitive advantage and innovation intelligibility & conditional & full mediation: value co-creation and the economic and relationship value it creates have no direct effect; dynamic and innovation capabilities carry the entire effect & Value co-creation in the AI innovation ecosystem has no significant direct effect on competitive advantage - the effect runs through dynamic and innovation capabilities (full mediation). \\
\addlinespace[2pt]
S11 & Ali et al. (2023) & Both & Perf: clinical, financial, organizational and technological outcomes (qualitative, financial gains quantified); CA: competitive position vs. rivals (patient choice, payer bargaining power, differentiation) & positive & prerequisites: qualified robotic and revenue-cycle team, cultural shift in the medical community, regulatory approval; barrier: high acquisition costs not covered by insurance & AI-enabled robotic surgery improves clinical, financial and technological outcomes and strengthens the competitive position of the healthcare organization. \\
\addlinespace[2pt]
S12 & Czarnitzki et al. (2023) & Perf. & Perf: firm productivity (sales- and value-added-based, production function) & positive &  & AI adoption is robustly associated with higher firm productivity in representative German data; the effect grows with usage intensity. \\
\addlinespace[2pt]
S13 & Huang \& Lee (2023) & Perf. & Perf: firm market value (stock-price reaction around the announcement) & positive & detailed announcements pay more overall but react with delay - vague ones only on the day itself; IT firms gain more, non-IT reactions not significant; negative wording correlates modestly negatively & Markets reward AI implementation announcements - more so when announcements are detailed (delayed but stronger reaction) and come from IT firms; non-IT reactions not significant. \\
\addlinespace[2pt]
S14 & Krakowski et al. (2023) & Both & Perf: competitive performance (game outcomes/ratings); CA: persistent performance differences (sustained advantage) across settings & conditional & substitution: traditional human capabilities become obsolete; complementation: new human-machine capabilities emerge that are unrelated or negatively related to traditional ones & AI adoption shifts the sources of competitive advantage: a new decision-making resource emerges at the human-AI intersection, unrelated to traditional capability. \\
\addlinespace[2pt]
S15 & Samadhiya et al. (2023) & Both & Perf: perceived firm performance (survey scale, expectation-worded items); CA: perceived sustainable competitiveness (survey scale; only one of five items competitor-referencing, rest mixes AIPRM embedding and environmental commitment) & positive & prerequisites: information and communication technology capability and technological readiness; firm fit & AI-based partner relationship management improves firm performance and sustainable competitiveness, provided information and communication technology capability and technological readiness are in place. \\
\addlinespace[2pt]
S16 & Wu \& Monfort (2023) & Perf. & Perf: perceived firm performance (survey scale) & positive & configurational: marketing capabilities, customer value co-creation, market orientation and AI strategy jointly form necessary and sufficient recipes for high performance & AI marketing strategy raises firm performance, but only as part of configurations with marketing capabilities, co-creation and market orientation - a configurational result. \\
\addlinespace[2pt]
S17 & Babina et al. (2024) & Perf. & Perf: growth in sales, employment, and market valuation & positive & channel: product innovation rather than process efficiency; gains concentrate among firms that were already large, which raises industry concentration (superstar dynamics) & AI-investing firms grow in sales, employment and value via product innovation; benefits concentrate among large firms and go along with rising industry concentration. \\
\addlinespace[2pt]
S18 & Cannas et al. (2024) & Perf. & Perf: case-reported competitiveness outcomes: costs, lead times, service level, quality, safety, sustainability (partly quantified: 30\% stock reduction, 4\% OEE gain) & positive & barriers as boundary conditions: data quality, AI skills, high investment needs, unclear economic benefits, lack of AI cost-analysis experience (9 barriers total in Table 6: financial, organisational, strategic, technological) & AI improves supply-chain competitiveness, mostly evidenced in production processes; the gains depend on overcoming data-quality, skill and investment-evaluation barriers. \\
\addlinespace[2pt]
S19 & Pesqueira et al. (2024) & Both & Perf: operational efficiency in tender management (processing times, decision quality, document accuracy), expert-rated; CA: expert-rated comparative competitiveness of the four companies in the tendering context & conditional & enablers: governance and ethical frameworks, regulatory compliance, organizational readiness, alignment with business goals; senior-management involvement identified as critical & AI adopters manage pharmaceutical tendering faster and with better decisions than non-adopters, but only sustainably within robust governance and compliance structures. \\
\addlinespace[2pt]
S20 & Pham et al. (2024) & Perf. & Perf: outpatient revenue, inpatient revenue, ROA, productivity, occupancy (ROA not significant in every specification) & mixed & larger-market-share hospitals are the adopters (self-selection, no test of who benefits more); without accounting for it the estimates mislead & Hospitals with larger market share adopt AI and gain in revenue, productivity and occupancy; naive comparisons that ignore which hospitals adopt AI misestimate the effects. \\
\addlinespace[2pt]
S21 & Sullivan \& Fosso (2024) & Perf. & Perf: perceived firm performance, process innovation, product innovation & conditional & central mediator: adaptive response to market changes; only AI analytics shows direct effects in a supplementary check; hostility weakens the automation path, strengthens the analytics path; only ten of eighteen conditional effects significant & AI capabilities pay off through the firm's adaptive response to market changes; environmental hostility/dynamism condition how much each AI capability contributes. \\
\addlinespace[2pt]
S22 & Sun et al. (2024) & Perf. & Perf: firm total factor productivity & positive & stronger for state-controlled, internationally oriented and innovative firms; channels: reduced information asymmetry (key), specialized division of labor and green innovation only potential; supply-chain digitalization policy amplifies; the effect weakens over time & AI raises firm productivity in China - most for state-controlled, international and innovative firms - but the productivity boost decays over time. \\
\addlinespace[2pt]
S23 & Zebec \& Indihar (2024) & Perf. & Perf: organizational performance (survey scale) & positive & mediators in sequence: decision-making quality and business-process performance carry the effect; process automation, organizational learning and process innovation as path-specific partial mediators & AI creates business value indirectly: through better decisions and business-process performance, complemented by automation, learning and process innovation. \\
\addlinespace[2pt]
S24 & Alam et al. (2025) & Both & Perf: value creation (survey scale); CA: sustainable competitive advantage (modeled as mediator) & positive & mediator: sustainable competitive advantage between AI adoption and value creation; dynamic capabilities have a negative direct effect on value creation (transition costs); technological turbulence moderates only the capabilities-to-value path, weakening it further & AI adoption creates value in hospitality directly and via sustainable competitive advantage; dynamic capabilities show a negative direct effect on value creation (transition costs), aggravated under technological turbulence. \\
\addlinespace[2pt]
S25 & Alwakid \& Dahri (2025) & Perf. & Perf: sustainable SME performance (survey scale) & positive & mediators: SME creativity and green innovation; green entrepreneurial orientation as parallel driver & AI capabilities raise sustainable SME performance through creativity and green innovation; AI infrastructure, proactive strategy and entrepreneurial risk-taking are the key resource inputs. \\
\addlinespace[2pt]
S26 & Banna \& Alam (2025) & Perf. & Perf: firm revenue growth & conditional & U-shape: AI hurts revenue below and helps above roughly USD 11.3M; works only when paired with own R\&D; only small and medium firms benefit & AI investment follows a U-shaped payoff: short-term revenue drag, long-term gains - and only firms coupling AI with R\&D strategy capture the upside. \\
\addlinespace[2pt]
S27 & Basnet et al. (2025) & Perf. & Perf: firm value (market valuation); R\&D, patents, and later productivity as channels & conditional & disclosure substance decides: actionable AI disclosures bring valuation gains, innovation, and later productivity; speculative or irrelevant ones bring nothing; silent or vague peers are penalized & Markets reward only substantive (actionable) AI disclosures - early adopters with real implementation plans gain value via innovation; vague AI talk earns nothing. \\
\addlinespace[2pt]
S28 & Bin-Nashwan \& Li (2025) & Perf. & Perf: perceived accounting-firm performance & conditional & mediators: green human capital and green structural capital carry the effect (no direct path modeled); green relational capital channel empty; moderator: sustainability culture & AI-infused knowledge improves accounting-firm performance through green human and structural capital, amplified by sustainability culture; relational capital plays no role. \\
\addlinespace[2pt]
S29 & Cao et al. (2025) & Perf. & Perf: efficiency and innovation outcomes (configurational) & conditional & configurational: no single AI function suffices; performance needs combinations, especially customer service and cybersecurity paired with other functions; emergent capabilities: adaptive learning, predictive analytics, uncertainty mitigation & Retail performance comes from combinations of AI-enabled functions, not from isolated AI uses; a directly configurational result. \\
\addlinespace[2pt]
S30 & Chakraborty (2025) & Both & Perf: firm/organisational performance (survey scales); CA: competitive advantage / competitive positioning (survey scale) & positive & enabler: technological readiness; moderator: leader narcissism shapes the adoption drivers and the effect of generative AI on competitive positioning & GenAI adoption improves hotels competitive advantage and performance, with leadership personality (narcissism) shaping how adoption translates into outcomes. \\
\addlinespace[2pt]
S31 & Chiu et al. (2025) & Perf. & Perf: profit, sustainability, and market-valuation efficiency & mixed & stage-specific: AI innovation lifts sustainability and market efficiency but not profit efficiency (size and leverage drive profit); firm scale raises profit efficiency yet lowers sustainability and market efficiency & AI-driven innovation lifts sustainability and market-valuation efficiency in semiconductors but not short-term profit efficiency - the payoff is stage-specific. \\
\addlinespace[2pt]
S32 & D’Amico et al. (2025) & Perf. & Perf: innovation performance (knowledge spillover of innovation) & conditional & complement: own R\&D investment (AI pays off only when paired with internal R\&D); AI substitutes for knowledge collaboration with certain external partners & AI adoption boosts innovation only when paired with the firms own R\&D investment - and partly replaces external knowledge collaboration. \\
\addlinespace[2pt]
S33 & Feng et al. (2025) & Perf. & Perf: innovation output and innovation efficiency & positive & stronger for large incumbents and growth-stage firms; executives with IT backgrounds amplify; highly innovative firms diversify supply chains to reduce resource risk & AI adoption raises innovation output and efficiency, most where firm maturity, scale and IT-savvy leadership provide absorptive conditions. \\
\addlinespace[2pt]
S34 & Fontanelli et al. (2025) & Perf. & Perf: volatility of productivity growth (not its level) & conditional & buyers of external AI see higher productivity volatility, in-house developers none; complementary human capital (share of technology engineers and technicians) softens the buyer effect & Off-the-shelf AI raises the riskiness of productivity outcomes; developing in-house or employing complementary technical staff neutralizes it. \\
\addlinespace[2pt]
S35 & Guo et al. (2025) & CA & CA: international advantage: product innovation and brand upgrading (global value-chain position) & positive & mediator: international marketing capabilities; platform embeddedness shows inverted-U effects (depth on product innovation, breadth on brand upgrading); medium embeddedness is optimal & AI adoption coupled with strategic agility upgrades SMEs global value-chain position through marketing capabilities - but only at moderate platform embeddedness (inverted-U). \\
\addlinespace[2pt]
S36 & Huang \& Lin (2025) & Perf. & Perf: financial ratios, productivity, market value (Tobin's q) & mixed & indicator split: financial performance and market value improve, productivity does not; first movers show partly negative or insignificant effects (no uniform first-mover advantage) & AI implementation lifts financial performance and market value, but first movers and top performers do not benefit uniformly across all indicators. \\
\addlinespace[2pt]
S37 & Kumar et al. (2025) & Perf. & Perf: perceived firm performance (survey scale) & positive & moderator: ethical leadership on the adoption-to-performance path; adoption drivers and barriers: uniqueness, information completeness, convenience, deceptiveness & GenAI adoption boosts perceived firm performance, more strongly under ethical leadership. \\
\addlinespace[2pt]
S38 & Mehta et al. (2025) & Both & Perf: firm performance (survey scale); CA: competitive advantage (modeled as mediator) & positive & manager personality traits (Big Five) drive adoption; competitive advantage mediates between adoption and performance; organizational culture moderates the effect of agreeableness on adoption & Individual personality traits shape AI adoption in key account management, which converts into firm performance through competitive advantage - culture conditions the path. \\
\addlinespace[2pt]
S39 & N. \& Roy (2025) & Perf. & Perf: firm performance (survey scale) & positive & mediator: marketing agility; moderator: technological turbulence amplifies the positive effects & AI-based stakeholder engagement improves performance through marketing agility - and the payoff is largest in technologically turbulent environments. \\
\addlinespace[2pt]
S40 & Sandeep et al. (2025) & CA & CA: perceived competitive advantage from AI-enabled recruitment (survey scale) & positive & HR competencies and open innovation build dynamic capabilities, which mediate; financial support and IT infrastructure act as enabling resources & AI adoption in recruitment yields competitive advantage only where dynamic capabilities, built from HR competencies and open innovation with financial/IT support, are in place. \\
\addlinespace[2pt]
S41 & Shi et al. (2025) & Perf. & Perf: enterprise value (Tobin's q) & positive & data assets mediate 45.6\% of the effect while a significant direct effect remains; managerial capability moderates; boundary conditions: no effect in heavily polluting industries unless the firm has a credible green transition, and none in asset-intensive firms; state-owned enterprises gain about twice as much as private ones & AI adoption raises enterprise value through data assets, and only under capable management and agile (non-asset-heavy) structures - the study explicitly maps when AI becomes a source of advantage. \\
\addlinespace[2pt]
S42 & Song et al. (2025) & Perf. & Perf: firm value (stock-price reaction) & negative & small firms hit harder; losses more negative in recent years; damage ranked by failure type: accuracy worst, then safety, privacy, fairness & AI service failures destroy market value - most for small firms and accuracy failures; the downside risk of AI implementation is real and priced. \\
\addlinespace[2pt]
S43 & Tao et al. (2025) & Perf. & Perf: new product performance (survey scale) & positive & mediator: AI capabilities (big-data analytics assimilation alone is insufficient); moderator: electronic supply-chain collaboration strengthens the indirect effect & Big-data/AI assimilation lifts new product performance only when converted into AI capabilities, and more so with electronic supply-chain collaboration. \\
\addlinespace[2pt]
S44 & Tingbani et al. (2025) & Perf. & Perf: firm growth (sales growth) & positive & labour productivity amplifies the growth effect; labour cost and labour share weaken it & AI investment raises firm growth modestly on average; labour-market conditions moderate the size of the gain - productivity-rich, labour-cost-light firms capture more of it. \\
\addlinespace[2pt]
S45 & Adams et al. (2026) & Perf. & Perf: growth in sales, employment, assets, markups; stock-return response to monetary policy & positive & adoption concentrated in larger, more productive firms; adopters more responsive to monetary policy surprises & Firms adopting AI-powered pricing grow faster in sales, employment and markups - AI pricing is a capability large productive firms exploit first. \\
\addlinespace[2pt]
S46 & Agag et al. (2026) & Perf. & Perf: firm value; wage inequality as mediating channel & positive & wage inequality partially mediates: AI human capital compresses inequality and thereby raises firm value; the effect is stronger in dynamic and complex industries; by sector, the value effect is strongest in travel and the inequality compression in hospitality & AI human capital raises firm value partly by compressing wage inequality - a workforce-centered pathway to AI-enabled advantage, strongest in dynamic sectors. \\
\addlinespace[2pt]
S47 & Alnofeli et al. (2026) & Both & Perf: profitability (survey scale); CA: competitive advantage (survey scale) & positive & mediator: marketing ambidexterity (balancing exploration and exploitation) carries the effect & AI-CRM capability lifts profitability and competitive advantage through marketing ambidexterity - capability without ambidexterity closes the value-realisation gap only partially. \\
\addlinespace[2pt]
S48 & Arshad et al. (2026) & Perf. & Perf: revenue growth; cost reduction & mixed & combining robotic process automation and AI adds revenue growth via price increases but no extra cost reduction (each technology alone does cut costs); a strategic revenue-growth intent amplifies the revenue gains & Combining robotic process automation and AI pays through higher revenue, not lower cost — and only for firms whose strategic intent targets growth. \\
\addlinespace[2pt]
S49 & Bai et al. (2026) & Perf. & Perf: stock price crash risk; firm value & mixed & crash risk rises through lower information transparency and managerial over-optimism (mechanisms); worse for low-transparency and resource-constrained firms; value gains may reflect short-term market optimism & Policy-driven AI investment inflates firm value but simultaneously raises crash risk - short-term valuation gains trade off against long-term stability. \\
\addlinespace[2pt]
S50 & Chen et al. (2026) & Perf. & Perf: investment efficiency; Tobin's q (long horizon) & positive & mechanisms: better sales forecasts, reporting quality, and product/process innovation; exposure to data-privacy regulation weakens the effect & AI adoption improves how firms allocate capital - unless regulatory friction (data-privacy exposure) blunts it; value shows up in Tobins q only over longer horizons. \\
\addlinespace[2pt]
S51 & Dinh (2026) & Perf. & Perf: new product performance (survey scale) & positive & direct effect plus a path through more radical digital product innovation (mediator); digital sustainability leadership acts as both mediator and moderator & AI knowledge capabilities lift new product performance in emerging-market SMEs, with digital sustainability leadership as the critical boundary condition. \\
\addlinespace[2pt]
S52 & Filieri et al. (2026) & Perf. & Perf: supply chain efficiency, sales performance, market value & conditional & inverted U-shape: AI benefits diminish at extreme tariff exposure (an optimal investment threshold); customer engagement needs strengthen the effect; tariff exposure itself induces AI integration & Trade tensions push firms into AI integration that improves efficiency, sales and value - up to a threshold where adversity overwhelms the compensation. \\
\addlinespace[2pt]
S53 & Giordino et al. (2026) & Perf. & Perf: ROA; Tobin's q; ESG pillar scores & positive &  & Organizational AI focus goes hand in hand with better financial performance and better environmental and social sustainability scores in European listed firms. \\
\addlinespace[2pt]
S54 & Kazakis (2026) & Perf. & Perf: firm efficiency & conditional & no unconditional effect; gains only with capable managers, stronger competitive pressure, stable institutional ownership, and access to long-term debt & AI investment alone does nothing for efficiency - gains materialize exclusively under managerial, competitive, ownership and financing conditions ("Conditional Gains"). \\
\addlinespace[2pt]
S55 & Li et al. (2026) & CA & CA: perceived competitive advantage (survey scale) & conditional & full mediation: no significant direct path; the effect runs entirely through distributive and procedural justice (fairness of outcomes and of process); supply chain complexity weakens the distributive channel but not the procedural one & Responsible AI builds competitive advantage through supply-chain justice - in complex networks, only procedural fairness remains a robust channel. \\
\addlinespace[2pt]
S56 & Li et al. (2026) & CA & CA: perceived sustained competitive advantage (survey scale) & positive & three GenAI affordances mediate between technological opportunism (a firm's readiness to seize new technology) and competitive advantage; competitive intensity amplifies only the creation-related path & A firm's readiness to seize new technology converts into competitive advantage through what GenAI enables; under intense competition, only the creative use of GenAI keeps paying. \\
\addlinespace[2pt]
S57 & Lin et al. (2026) & Perf. & Perf: total factor productivity & positive & internal control quality partially mediates: AI improves governance and risk management, which in turn raises total factor productivity & AI raises firm productivity partly by improving internal control quality - governance is a transmission channel, not just a moderator. \\
\addlinespace[2pt]
S58 & Lin et al. (2026) & Perf. & Perf: investment efficiency; firm value & positive & stronger in knowledge-intensive and R\&D-intensive firms and firms with sound internal controls; mechanisms: innovation capability and financial health & AI enhances investment efficiency and firm value most where knowledge intensity, R\&D and internal controls give the firm the capacity to absorb it. \\
\addlinespace[2pt]
S59 & Liu et al. (2026) & Perf. & Perf: likelihood of M\&A; long-term firm value & positive & stronger with abundant cash holdings, less-developed market environments, private ownership & AI adoption makes firms more acquisitive and raises long-term value - internal resources and institutional context set the boundary. \\
\addlinespace[2pt]
S60 & Liu et al. (2026) & Perf. & Perf: productivity gains + firm growth & positive & regulatory clarity and supportive enforcement strengthen the payoff; confirmed causally in an experiment & AI-enabled innovation converts into productivity and growth only under clear, supportive regulation - institutional environment is the amplifier. \\
\addlinespace[2pt]
S61 & Mukherjee et al. (2026) & Perf. & Perf: service performance (survey scale) & positive &  & AI adoption intention lifts hotel service performance; drivers: ease of use, competitive pressure, policy support - the infrastructure estimate is negative. \\
\addlinespace[2pt]
S62 & Patel (2026) & Perf. & Perf: stock market reaction to patent grants & negative & firm-specific stock volatility softens the negative reaction (moderator); asset tangibility does not; the effect appears only in the 2010s and 2020s & Markets react negatively to AI-heavy service patents in manufacturing — a caution against assuming AI-based service moves create short-term market value. \\
\addlinespace[2pt]
S63 & Renfei et al. (2026) & Perf. & Perf: corporate sustainability performance including an economic dimension (cost reduction, revenue from waste, spin-offs) & positive & prescriptive (action-recommending) capabilities fail as part of the AI capability bundle; firm type is a boundary condition: strong effect for Industry 4.0 (digitally advanced) manufacturers, much weaker but not absent for traditional firms & AI capabilities raise sustainability and economic performance mainly in digitally advanced (Industry 4.0) manufacturers; action-recommending (prescriptive) AI fails, and traditional firms barely benefit. \\
\addlinespace[2pt]
S64 & Singh et al. (2026) & Perf. & Perf: ROA; firm value & conditional & small positive ROA effect, stronger for R\&D-intensive firms; firm-value gains only for firms with good sustainability records; stronger under Democratic administrations (political environment as condition) & AI innovation pays modestly and unevenly: R\&D intensity, sustainability credentials and even the political climate condition whether markets reward it. \\
\addlinespace[2pt]
S65 & Song et al. (2026) & Perf. & Perf: market reactions; long-term operating performance & negative & time path: short-lived positive investor reaction, then negative long-run stock returns and sustained underperformance; the gap between claims and patents is what drives the effect & AI talk without AI substance backfires: markets initially reward the narrative, then punish rhetoric that outpaces implementation. \\
\addlinespace[2pt]
S66 & Ullah et al. (2026) & Perf. & Perf: organizational performance (survey scale) & positive & team dynamics, corporate innovation capabilities, high-commitment workforce (mediators) & AI adoption converts into performance where cohesive teams, innovation capability and a committed workforce absorb it - human/organizational enablers unlock the value. \\
\addlinespace[2pt]
S67 & Wang \& Zhang (2026) & Perf. & Perf: entrepreneurial success (survey scale) & conditional & full mediation: strategic resilience fully carries the effect of AI capability on venture success; configurational analysis finds several equally effective condition combinations for high performance & AI capability makes ventures successful only via organizational resilience - the mechanism is organizational, not technological; multiple configurations lead there. \\
\addlinespace[2pt]
\end{longtable}
\end{landscape}
\endgroup
